\chapter{总结与展望}


\section{论文总结}
本文围绕资源受限微型无人机在复杂未知环境中的协同感知与自主导航问题，针对集群测距、协同定位以及安全导航等关键技术展开了系统性研究，并完成了从算法设计到原型系统实现的完整验证。主要工作总结如下：

首先，针对微型无人机集群在高密度与高动态环境中测距不稳定的问题，本文在第三章提出了一种面向相对定位的鲁棒集群测距协议。通过对随机测距周期进行理论建模与参数优化，有效缓解了测距冲突与报文不匹配问题；同时引入不完全报文匹配下的补偿测距机制，提高了测距成功率。在此基础上，结合期望最大化算法，提出了一种自适应扩展卡尔曼滤波方法，实现了过程噪声与测量噪声的在线估计，从而在缺乏先验噪声信息的条件下提升了协同定位的精度与鲁棒性。

其次，针对无人机在复杂环境中自主导航面临的安全性不足问题，第四章提出了一种风险感知强化学习方法。该方法通过对回报分布进行建模，引入风险评估机制，使无人机在决策过程中能够兼顾收益与潜在风险，从而有效降低低概率高风险事件带来的影响。同时，通过对模型进行量化与压缩，实现了算法在资源受限嵌入式平台上的高效部署，满足了实时导航的需求。

最后，在第五章中，本文将所提出的协同定位方法与自主导航算法进行系统集成，构建了一个完整的微型无人机集群自主导航原型系统。该系统基于Crazyflie 2.1平台，采用“感知—定位—决策—控制”的分层架构，实现了多无人机之间的可靠测距、协同定位以及领航-跟随模式下的自主导航与避障。通过真实飞行实验验证，结果表明所构建系统能够在资源受限条件下稳定运行，具备良好的实时性与鲁棒性，验证了本文方法的工程可行性。

综上所述，本文从理论方法、算法设计到系统实现，系统性地解决了微型无人机集群在复杂环境中协同定位与安全导航的关键问题，为其在实际场景中的应用提供了重要参考。

\section{后续研究方向}
在前文工作的基础上，尽管本文已在微型无人机集群协同定位与自主导航方面取得了一定进展，但受限于研究条件与系统复杂性，相关方法在规模扩展、环境适应性以及协同智能水平等方面仍存在进一步提升空间。随着无人机技术与人工智能方法的不断发展，面向复杂动态环境的高效、安全与智能化集群系统将成为重要研究方向。基于此，本文结合第三章至第五章的研究内容，从协同测距与定位方法、风险感知决策机制以及系统集成与多机协同等方面，对未来可能的研究方向进行如下展望。
\subsection{集群测距与自适应定位方法的进一步研究（第三章）}
第三章主要解决了集群测距冲突与噪声未知条件下的协同定位问题，但仍存在一定的提升空间：

首先，在测距协议方面，本文采用随机周期与补偿机制来缓解冲突，但在更高密度集群或强干扰环境中，测距成功率仍可能下降。未来可研究基于时隙调度或自组织通信机制的测距协议，实现更加稳定和可预测的通信性能。同时，可结合图论方法，对测距拓扑结构进行优化，提高整体网络连通性。

其次，在测距误差建模方面，当前方法主要通过滤波进行补偿，缺乏对误差来源的精细建模。未来可针对非视距误差、多径效应等问题，引入环境感知或数据驱动的误差建模方法，从源头上提升测距精度。

再次，在自适应滤波方法方面，本文基于期望最大化算法实现了噪声协方差的在线估计，但其收敛速度和稳定性仍依赖一定条件。未来可研究更高阶自适应滤波方法（如UKF、IEKF或基于学习的滤波方法），并探索参数学习与状态估计的联合优化机制，以提升在强非线性场景下的性能。

最后，在分布式协同定位方面，当前方法仍存在一定集中处理特征。未来可进一步发展完全分布式定位算法，减少对单节点的依赖，提高系统的鲁棒性与可扩展性。

\subsection{风险感知强化学习导航方法的深化研究（第四章）}

第四章基于CVaR风险度量构建了风险感知强化学习导航方法，在提升无人机规避高风险行为方面取得了一定效果。然而，在复杂动态环境及工程落地过程中，仍存在进一步优化空间，具体体现在以下几个方面：

首先，在模型泛化能力方面，当前方法主要依赖于特定环境下的训练数据，其策略性能对环境分布具有一定依赖性。当环境发生变化（如障碍物分布、传感器噪声特性或动力学参数变化）时，策略可能出现性能退化甚至失效。为此，未来可引入域随机化（Domain Randomization）方法，通过在训练阶段对环境参数进行大范围扰动，使模型学习到更具鲁棒性的策略表示。此外，可结合迁移学习（Transfer Learning）或元学习（Meta-learning）方法，实现策略在不同任务或环境之间的快速适配，从而提升系统在未知场景中的泛化能力与实用性。

其次，在安全约束机制方面，尽管CVaR从风险分布角度对策略进行了约束，但其本质仍属于“软约束”，难以提供严格的安全保证。在实际无人机系统中，碰撞避免等安全约束往往具有“硬约束”特性。未来可考虑将强化学习与传统控制方法相结合，例如引入模型预测控制（Model Predictive Control, MPC）框架，在优化过程中显式加入状态与控制约束；或利用控制屏障函数（Control Barrier Function, CBF）构建安全过滤层，对强化学习输出的控制指令进行实时修正，从而形成“高层学习决策 + 低层安全约束”的分层控制结构，在保证安全性的同时兼顾策略的灵活性与最优性。

再次，在动态环境适应能力方面，当前方法主要针对静态或弱动态环境进行建模，对于快速变化的环境（如移动障碍物或多无人机交互场景）仍缺乏有效应对机制。未来可进一步研究时序建模方法（如引入循环神经网络RNN或注意力机制），增强策略对环境动态变化的感知能力；

此外，在轻量化与嵌入式部署方面，虽然本文已通过模型压缩降低了计算开销，但在更严格的资源受限平台（如纳型无人机）上仍存在实时性与功耗瓶颈。未来可从算法与硬件协同优化角度出发，进一步研究神经网络剪枝（Pruning）、知识蒸馏（Knowledge Distillation）以及量化感知训练（Quantization-Aware Training）等方法，在保证性能的前提下降低模型复杂度。同时，可探索基于专用加速硬件（如FPGA或边缘AI芯片）的部署方案，以实现高效、低功耗的实时推理。


\subsection{系统集成与多机协同能力提升（第五章）}

第五章完成了协同定位与风险感知导航方法的系统集成，构建了基于微型无人机平台的自主导航原型系统，并通过实验验证了系统的可行性与有效性。然而，从实际应用角度出发，系统在规模扩展、环境适应性以及协同智能水平等方面仍有进一步提升空间，未来可从以下几个方向展开研究：

首先，在系统规模扩展方面，当前系统主要面向小规模无人机集群进行设计与验证。当集群规模进一步扩大时，通信冲突、信息延迟及系统稳定性问题将更加突出。未来可研究面向大规模集群的分层式系统架构与通信调度机制，通过引入局部协同与全局协调相结合的策略，提高系统的可扩展性与运行效率。

其次，在复杂环境适应性方面，现有系统主要在室内或受控环境中进行验证，对于室外复杂环境（如动态障碍物、光照变化及弱定位条件）仍缺乏充分测试。未来可结合多传感器融合感知与环境建图技术，提升无人机在未知环境中的感知与定位能力，从而增强系统在真实场景中的适用性与鲁棒性。

再次，在多机协同决策能力方面，当前系统主要基于领航–跟随模式实现基础协同，其协同策略相对固定，难以适应复杂动态任务需求。未来可进一步引入多智能体强化学习（Multi-Agent Reinforcement Learning, MARL）方法，通过多无人机之间的协同学习与信息交互，实现去中心化的协同决策与自主任务分配。该方法能够使各无人机在局部观测条件下学习到协同避障与路径规划策略，从而在多障碍物、动态干扰及多机交互场景中显著提升系统整体性能与自适应能力。同时，可结合“集中训练–分布执行”机制，在保证实时性的前提下提高系统的可扩展性与稳定性。

最后，在软硬件协同优化方面，随着任务复杂度的提升，对计算资源与能耗提出了更高要求。未来可从系统层面出发，对计算、通信与控制模块进行联合优化，并结合嵌入式硬件加速技术，实现更高效、低功耗的无人机集群系统设计。